\section{Uso de modelos preentrenados}\label{sec:modelos}

El proyecto no entrena ni ajusta modelos: usa modelos fundacionales a través de API y concentra el
esfuerzo en los datos, la recuperación y la orquestación. Es una decisión, no una carencia. El
conocimiento de una ciudad cambia (cierres, horarios, nuevos locales), y con RAG se actualiza
reconstruyendo el corpus en minutos, con la fuente y la licencia trazables. Un ajuste fino
(\emph{fine-tuning}), en cambio, congelaría ese conocimiento en los pesos, no permitiría citar de
dónde sale cada lugar y exigiría un conjunto de entrenamiento que no existe. La
Tabla~\ref{tab:modelos} resume, al estilo de una ficha de modelo
(\emph{model card})~\cite{mitchell2019modelcards}, los modelos empleados.

\begin{table}[h]
\centering\scriptsize
\caption{Modelos utilizados y su papel.}\label{tab:modelos}
\begin{tabularx}{\columnwidth}{@{}L{2.15cm}Y@{}}
\toprule
\textbf{Modelo} & \textbf{Uso y configuración} \\
\midrule
Claude Haiku 4.5 (Bedrock, perfil de inferencia UE) & Agente, extracción del resumen y elecciones estructuradas. Temperatura 0,7; máximo de 4.096 tokens de salida. Elegido por admitir herramientas y salida estructurada nativa con JSON Schema, por su latencia y por su precio (1 y 5\,\$ por millón de tokens de entrada y salida) \\
Titan Text Embeddings V2 (Bedrock, misma región) & Vectores de 1.024 dimensiones normalizados para documentos y consultas; multilingüe: una pregunta en español encuentra texto en inglés \\
Bedrock Guardrails & Clasificadores gestionados de ataques de \emph{prompt}, contenido dañino, datos de pago y temas denegados \\
Amazon Nova Pro (Bedrock) & Juez de la evaluación, temperatura 0; de otra familia de modelos, para evitar el sesgo de autoevaluación \\
Nemotron 3 Super 120B (NVIDIA) & Alternativa para desarrollo local y referencia de la evaluación del agente \\
\bottomrule
\end{tabularx}
\end{table}

\textbf{Selección del modelo del agente.} Haiku~4.5 es el modelo de Bedrock desde la integración
con este servicio (septiembre), por precio, latencia y soporte de herramientas y JSON Schema nativo,
y se sirve desde la UE (RGPD). La evaluación de conversaciones (§\ref{sec:evalagente}) confirmó
después la elección (Tabla~\ref{tab:modelcomp}): iguala a Nemotron dentro de la tolerancia (33,3
frente a 32,7 escenarios) con un 17\,\% menos de tokens y un 13\,\% menos de llamadas. Un modelo
mayor multiplicaría el coste por turno para ganar, como mucho, los tres escenarios que aún fallan
(ediciones en español cuya segunda mitad queda sin hacer), algo que se aborda antes en el
\emph{prompt}. GLM-5.3-Flash quedó descartado porque su prueba de herramientas agotaba siempre el
tiempo de espera, aunque sigue siendo útil en local
para los pasos estructurados. La temperatura de 0,7 da una prosa natural y variada; la forma de
la respuesta no depende de ella, porque la fijan los esquemas y las herramientas. El juez, en
cambio, usa temperatura 0 para ser reproducible.

\textbf{Selección del modelo de \emph{embeddings}.} Titan Text Embeddings V2 es multilingüe (una
pregunta en español encuentra texto en inglés sin traducir), se ejecuta en la misma región que los
datos y cuesta 0,02\,\$ por millón de tokens, de modo que indexar una ciudad cuesta céntimos.

\begin{table}[h]
\centering\scriptsize
\caption{Modelos evaluados para el agente (evaluación de conversaciones, 36 escenarios $\times$3).}\label{tab:modelcomp}
\begin{tabularx}{\columnwidth}{@{}Yccc@{}}
\toprule
& \textbf{Haiku 4.5} & \textbf{Nemotron 3 Super} & \textbf{GLM-5.3-Flash} \\
\midrule
Proveedor & Bedrock (UE) & NVIDIA & NVIDIA \\
Uso de herramientas & sí & sí & no \\
JSON Schema nativo & sí & no & sí \\
Escenarios superados & 92,6\,\% & 90,7\,\% & -- \\
Comprobaciones superadas & 577/585 & 575/585 & -- \\
Tokens de entrada & 0,99 M & 1,20 M & -- \\
Llamadas al modelo & 295 & 338 & -- \\
Coste de la evaluación & 1,15\,\$ & sin precio público & -- \\
\bottomrule
\end{tabularx}
\end{table}

El acceso a los modelos pasa por un puerto (\code{LLMProvider}) con adaptadores para Bedrock y
para la API compatible con OpenAI de NVIDIA, por lo que cambiar de modelo es una cuestión de configuración. Las
salidas estructuradas usan el modo nativo del proveedor cuando el modelo está verificado
y, si no, el esquema como instrucción con validación posterior; ante una salida inválida
no se reintenta a ciegas, sino que se recurre a una alternativa determinista (los mejores
candidatos de la búsqueda). Cada \emph{prompt} se versiona con el \emph{hash} de su plantilla, que se
guarda en las trazas y en los resultados de evaluación para no comparar ejecuciones con
\emph{prompts} distintos.
